% ICPC Team Reference Document (TRD) template
% Compile with pdfLaTeX (twice if references are added).
% Before printing, replace the team metadata below and check the final PDF page count.

\documentclass[10pt,a4paper]{article}

% ---------- Encoding and fonts ----------
\usepackage[T2A]{fontenc}
\usepackage[utf8]{inputenc}
\usepackage[english,russian]{babel}
\usepackage{lmodern}
\usepackage{microtype}

% ---------- Page layout ----------
\usepackage[
a4paper,
left=12mm,
right=12mm,
top=17mm,
bottom=12mm,
headheight=14pt,
headsep=4mm
]{geometry}
\usepackage{fancyhdr}
\usepackage{multicol}
\setlength{\columnsep}{7mm}
\setlength{\columnseprule}{0pt}
\setlength{\parindent}{0pt}
\setlength{\parskip}{2pt plus 1pt}
\raggedcolumns

% ---------- Mathematics, tables, lists ----------
\usepackage{amsmath,amssymb,mathtools}
\usepackage{booktabs,array}
\usepackage{enumitem}
\setlist{nosep,leftmargin=*}

% ---------- Code ----------
\usepackage{xcolor}
\usepackage{listings}
\lstdefinestyle{cpp}{
	language=C++,
	basicstyle=\ttfamily\footnotesize,
	keywordstyle=\bfseries,
	commentstyle=\itshape,
	stringstyle=\ttfamily,
	numbers=none,
	tabsize=4,
	keepspaces=true,
	showstringspaces=false,
	columns=fullflexible,
	breaklines=true,
	breakatwhitespace=false,
	frame=single,
	framerule=0.25pt,
	framesep=3pt,
	aboveskip=3pt,
	belowskip=3pt
}
\lstset{style=cpp}

% ---------- Compact headings ----------
\usepackage{titlesec}
\titleformat{\section}{\large\bfseries}{\thesection.}{0.45em}{}
\titleformat{\subsection}{\normalsize\bfseries}{\thesubsection}{0.4em}{}
\titlespacing*{\section}{0pt}{5pt plus 2pt minus 1pt}{2pt}
\titlespacing*{\subsection}{0pt}{4pt plus 1pt minus 1pt}{1pt}

% ---------- Team metadata: EDIT THESE ----------
\newcommand{\UniversityName}{Yaroslavl State University}
% Sort surnames alphabetically, as required by the regional TRD rules.
\newcommand{\TeamSurnames}{Chirkov, Rakushin, Vorobyov}
\newcommand{\ContestName}{Central Russia Regional Contest}

% ---------- Header: required metadata and page number ----------
\pagestyle{fancy}
\fancyhf{}
\fancyhead[L]{\footnotesize \UniversityName\ (\TeamSurnames)}
\fancyhead[R]{\footnotesize \thepage}
\renewcommand{\headrulewidth}{0.25pt}
\renewcommand{\footrulewidth}{0pt}
\fancypagestyle{plain}{%
	\fancyhf{}
	\fancyhead[L]{\footnotesize \UniversityName\ (\TeamSurnames)}
	\fancyhead[R]{\footnotesize \thepage}
	\renewcommand{\headrulewidth}{0.25pt}
}

% ---------- Convenience commands ----------
\newcommand{\complexity}[1]{\textit{Complexity:}~#1}
\newcommand{\invariant}{\textit{Invariant:}~}
\newcommand{\pitfall}{\textit{Pitfall:}~}
\newcommand{\code}[1]{\texttt{#1}}

% Count physical pages shipped to the PDF and warn if the 25-page limit is exceeded.
\newcounter{trdpagecount}
\AddToHook{shipout/after}{\stepcounter{trdpagecount}}
\AtEndDocument{%
	\clearpage
	\ifnum\value{trdpagecount}>25
	\PackageWarningNoLine{TRD}{PDF has \arabic{trdpagecount} pages; regional limit is 25 single-sided A4 pages}
	\fi
}

\begin{document}

\subsection{FFT}

Коэффициенты многочлена хранятся в порядке возрастания
степеней: $P[i]$ — коэффициент при $x^i$.

\complexity{Время: $O(n \log n)$, память: $O(n)$.}

\begin{lstlisting}
	struct cmpl{
		double a, b;
		cmpl(double a = 0, double b = 0): a(a), b(b) {}
		
		cmpl operator+(const cmpl &other) const{
			return cmpl(a + other.a, b + other.b);
		}
		
		cmpl operator-(const cmpl &other) const{
			return cmpl(a - other.a, b - other.b);
		}
		
		cmpl operator*(const cmpl &other) const{
			return cmpl(a * other.a - b * other.b,
			a * other.b + b * other.a);
		}
		
		cmpl operator/(const double c) const{
			return cmpl(a / c, b / c);
		}
		
		double &real(){
			return a;
		}
		
		double &imag(){
			return b;
		}
	};
	
	const double PI = atan2(0, -1);
	vector<vector<cmpl>> w;
	
	void precompute_w(int max_level){
		w.resize(max_level + 1);
		w[0].resize(1, 1);
		
		for (int l = 1; l <= max_level; l++){
			int n = (1 << l);
			w[l].resize(n / 2);
			
			cmpl firstw(cos(2 * PI / n), sin(2 * PI / n));
			
			for (int j = 0; j < n / 2; j++){
				if (j % 2 == 0){
					w[l][j] = w[l - 1][j / 2];
				}
				else{
					w[l][j] = w[l - 1][j / 2] * firstw;
				}
			}
		}
	}
	
	void fft(vector<cmpl> &P, bool inv = false){
		int n = P.size();
		int logn = 0;
		
		while ((1 << logn) != n){
			logn++;
		}
		
		vector<int> bit_rev(n);
		bit_rev[0] = 0;
		
		for (int j = 1; j < n; j++){
			bit_rev[j] = bit_rev[j / 2] / 2
			+ ((j & 1) << (logn - 1));
		}
		
		for (int j = 0; j < n; j++){
			if (j < bit_rev[j]){
				swap(P[bit_rev[j]], P[j]);
			}
		}
		
		for (int lvl = 1; lvl <= logn; lvl++){
			int len = 1 << lvl;
			int half = len >> 1;
			
			for (int st = 0; st < n; st += len){
				for (int j = 0; j < half; j++){
					cmpl tw = w[lvl][j];
					
					if (inv) tw.imag() = -tw.imag();
					
					cmpl u = P[st + j];
					cmpl v = tw * P[st + j + half];
					
					P[st + j] = u + v;
					P[st + j + half] = u - v;
				}
			}
		}
		
		if (inv){
			for (int i = 0; i < n; i++){
				P[i] = P[i] / n;
			}
		}
	}
	
	vector<ll> multiply(const vector<ll>& P,
	const vector<ll>& Q) {
		if (P.empty() || Q.empty()) return {};
		
		int need = (int)P.size() + (int)Q.size() - 1;
		int n = 1, logn = 0;
		
		while (n < need) {
			n *= 2;
			logn++;
		}
		
		precompute_w(logn);
		
		vector<cmpl> A(n);
		
		for (int j = 0; j < n; j++) {
			A[j] = cmpl(
			j < P.size() ? P[j] : 0,
			j < Q.size() ? Q[j] : 0
			);
		}
		
		fft(A);
		
		for (int j = 0; j < n; j++) {
			A[j] = A[j] * A[j];
		}
		
		fft(A, true);
		
		vector<ll> R(need);
		
		for (int j = 0; j < need; j++) {
			R[j] = llround(A[j].imag() / 2.0);
		}
		
		return R;
	}
\end{lstlisting}

\subsection{Dinic}

\begin{lstlisting}[style=cpp]
	#include <algorithm>
	#include <cassert>
	#include <limits>
	#include <queue>
	#include <vector>
	using namespace std;
	
	template<typename T>
	class max_flow {
		public:
		struct edge {
			int to, rev;
			T capacity, flow;
		};
		
		struct path {
			T flow;
			vector<int> path, edges;
		};
		
		private:
		int n;
		vector<vector<edge>> g;
		vector<int> dist, head;
		int source = -1, sink = -1;
		
		bool bfs() {
			dist.assign(n, n);
			dist[source] = 0;
			vector<int> que{source};
			
			for (int ptr = 0; ptr < (int)que.size(); ptr++) {
				int v = que[ptr];
				
				for (const auto &[u, rev, capacity, flow] : g[v]) {
					if (dist[u] > dist[v] + 1 && flow < capacity) {
						dist[u] = dist[v] + 1;
						que.push_back(u);
					}
				}
			}
			
			return dist[sink] != n;
		}
		
		T dfs(int v, T f) {
			if (v == sink)
			return f;
			
			for (; head[v] < (int)g[v].size(); head[v]++) {
				auto &[u, rev, capacity, flow] = g[v][head[v]];
				
				if (dist[u] == dist[v] + 1 && flow < capacity) {
					if (T delta = dfs(u, min(f, capacity - flow))) {
						flow += delta;
						g[u][rev].flow -= delta;
						return delta;
					}
				}
			}
			
			return 0;
		}
		
		// Requires the positive-flow graph to be acyclic.
		// Can recurse infinitely if there are positive-flow cycles.
		T decompose_dfs_acyclic(
		int v, T f,
		vector<int> &vertices,
		vector<int> &edge_ids
		) {
			if (v == sink) {
				vertices.push_back(sink);
				return f;
			}
			
			for (; head[v] < (int)g[v].size(); head[v]++) {
				auto &[u, rev, capacity, flow] = g[v][head[v]];
				
				if (flow > 0) {
					if (T delta = decompose_dfs_acyclic(
					u, min(f, flow), vertices, edge_ids
					)) {
						flow -= delta;
						g[u][rev].flow += delta;
						
						vertices.push_back(v);
						edge_ids.push_back(head[v]);
						return delta;
					}
				}
			}
			
			return 0;
		}
		
		public:
		max_flow(int n = 0) : n(n), g(n) {}
		
		int size() const {
			return n;
		}
		
		// Returns an index of the added edge in g[from].
		int add(
		int from, int to,
		T forward_capacity,
		T backward_capacity = 0
		) {
			assert(0 <= from && from < n);
			assert(0 <= to && to < n);
			assert(forward_capacity >= 0);
			assert(backward_capacity >= 0);
			
			int id = (int)g[from].size();
			
			g[from].push_back({
				to, (int)g[to].size() + (from == to),
				forward_capacity, 0
			});
			
			g[to].push_back({
				from, id, backward_capacity, 0
			});
			
			return id;
		}
		
		// Returns i-th edge of vertex v.
		edge get_edge(int v, int id) const {
			return g[v][id];
		}
		
		void clear_flow() {
			for (int v = 0; v < n; v++) {
				for (auto &e : g[v]) {
					e.flow = 0;
				}
			}
		}
		
		// U must be able to store the total maximum flow.
		// If source or sink changes, call clear_flow() first.
		template<typename U>
		U solve(int source_, int sink_) {
			assert(0 <= source_ && source_ < n);
			assert(0 <= sink_ && sink_ < n);
			assert(source_ != sink_);
			
			source = source_;
			sink = sink_;
			
			U flow = 0;
			
			while (bfs()) {
				head.assign(n, 0);
				
				while (T delta = dfs(
				source, numeric_limits<T>::max()
				)) {
					flow += delta;
				}
			}
			
			return flow;
		}
		
		// Returns true if v belongs to the source side of mincut.
		// Requires solve() to be called before.
		bool left_of_mincut(int v) const {
			return dist[v] != n;
		}
		
		// Fast decomposition.
		// Requires the positive-flow graph to be acyclic.
		vector<path> decompose_dfs_acyclic() {
			assert(source != -1 && sink != -1);
			
			vector<path> paths;
			vector<int> vertices, edge_ids;
			
			head.assign(n, 0);
			
			while (T flow = decompose_dfs_acyclic(
			source,
			numeric_limits<T>::max(),
			vertices,
			edge_ids
			)) {
				reverse(vertices.begin(), vertices.end());
				reverse(edge_ids.begin(), edge_ids.end());
				
				paths.push_back({flow, vertices, edge_ids});
				
				vertices.clear();
				edge_ids.clear();
			}
			
			return paths;
		}
		
		// General decomposition.
		// Works even if the positive-flow graph contains cycles.
		vector<path> decompose_bfs_general() {
			assert(source != -1 && sink != -1);
			
			vector<path> paths;
			
			while (true) {
				vector<int> pv(n, -1), pe(n, -1);
				queue<int> q;
				
				pv[source] = source;
				q.push(source);
				
				// Find a path using only positive-flow edges.
				while (!q.empty() && pv[sink] == -1) {
					int v = q.front();
					q.pop();
					
					for (int id = 0; id < (int)g[v].size(); id++) {
						const edge &e = g[v][id];
						
						if (e.flow <= 0 || pv[e.to] != -1)
						continue;
						
						pv[e.to] = v;
						pe[e.to] = id;
						q.push(e.to);
					}
				}
				
				if (pv[sink] == -1)
				break;
				
				// Find the bottleneck.
				T pushed = numeric_limits<T>::max();
				
				for (int v = sink; v != source; v = pv[v]) {
					pushed = min(
					pushed,
					g[pv[v]][pe[v]].flow
					);
				}
				
				path cur;
				cur.flow = pushed;
				
				// Subtract the extracted path from the flow.
				for (int v = sink; v != source; v = pv[v]) {
					int p = pv[v];
					int id = pe[v];
					
					int to = g[p][id].to;
					int rev = g[p][id].rev;
					
					g[p][id].flow -= pushed;
					g[to][rev].flow += pushed;
					
					cur.path.push_back(v);
					cur.edges.push_back(id);
				}
				
				cur.path.push_back(source);
				
				reverse(cur.path.begin(), cur.path.end());
				reverse(cur.edges.begin(), cur.edges.end());
				
				paths.push_back(cur);
			}
			
			return paths;
		}
	};
\end{lstlisting}

\subsection*{Suffix Array}

\begin{lstlisting}[style=cpp]
	struct suffix_array {
		string s;
		int n;
		vector<int> p, c;
		
		suffix_array(string s = "") {
			build(s);
		}
		
		void radix_sort(vector<pair<pair<int, int>, int>> &cc) {
			int n = cc.size();
			
			{
				vector<int> cnt(n);
				
				for (auto x : cc) {
					cnt[x.first.second]++;
				}
				
				vector<pair<pair<int, int>, int>> cc_nxt(n);
				vector<int> pos(n);
				
				pos[0] = 0;
				for (int i = 1; i < n; i++) {
					pos[i] = pos[i - 1] + cnt[i - 1];
				}
				
				for (auto x : cc) {
					int i = x.first.second;
					cc_nxt[pos[i]] = x;
					pos[i]++;
				}
				
				cc = cc_nxt;
			}
			
			{
				vector<int> cnt(n);
				
				for (auto x : cc) {
					cnt[x.first.first]++;
				}
				
				vector<pair<pair<int, int>, int>> cc_nxt(n);
				vector<int> pos(n);
				
				pos[0] = 0;
				for (int i = 1; i < n; i++) {
					pos[i] = pos[i - 1] + cnt[i - 1];
				}
				
				for (auto x : cc) {
					int i = x.first.first;
					cc_nxt[pos[i]] = x;
					pos[i]++;
				}
				
				cc = cc_nxt;
			}
		}
		
		void build(string str) {
			s = str;
			s += "$";
			
			n = s.size();
			
			p.resize(n);
			c.resize(n);
			
			vector<pair<char, int>> temp(n);
			
			for (int i = 0; i < n; i++) {
				temp[i] = {s[i], i};
			}
			
			sort(temp.begin(), temp.end());
			
			p[0] = temp[0].second;
			
			for (int i = 1; i < n; i++) {
				p[i] = temp[i].second;
				
				if (temp[i].first == temp[i - 1].first) {
					c[temp[i].second] = c[temp[i - 1].second];
				}
				else {
					c[temp[i].second] = c[temp[i - 1].second] + 1;
				}
			}
			
			vector<pair<pair<int, int>, int>> cc(n);
			
			int k = 0;
			
			while ((1 << k) < n) {
				for (int i = 0; i < n; i++) {
					cc[i] = {{c[i], c[(i + (1 << k)) % n]}, i};
				}
				
				radix_sort(cc);
				
				p[0] = cc[0].second;
				
				for (int i = 1; i < n; i++) {
					p[i] = cc[i].second;
					
					if (cc[i].first == cc[i - 1].first) {
						c[cc[i].second] = c[cc[i - 1].second];
					}
					else {
						c[cc[i].second] = c[cc[i - 1].second] + 1;
					}
				}
				
				k++;
			}
		}
		
		vector<int> get_sa() const {
			return p;
		}
	};
\end{lstlisting}

\complexity{Время: $O(n\log n)$, память: $O(n)$.}

\pitfall Символ \texttt{\$} должен отсутствовать в исходной строке и быть лексикографически меньше всех её символов.

\subsection*{Extended Euclidean Algorithm}

\begin{lstlisting}[style=cpp]
	ll gcd(ll a, ll b, ll &x, ll &y) {
		if (a == 0) {
			x = 0;
			y = 1;
			return b;
		}
		
		ll x1, y1;
		ll d = gcd(b % a, a, x1, y1);
		
		x = y1 - (b / a) * x1;
		y = x1;
		
		return d;
	}
\end{lstlisting}

\complexity{Time: $O(\log \min(a,b))$, memory: $O(\log \min(a,b))$.}
\subsection*{Z-function}

\begin{lstlisting}[style=cpp]
	vector<int> z_func(string &s) {
		int n = s.size();
		vector<int> z(n, 0);
		
		for (int i = 1, l = 0, r = 0; i < n; i++) {
			if (i <= r) {
				z[i] = min(r - i + 1, z[i - l]);
			}
			
			while (i + z[i] < n &&
			s[z[i]] == s[i + z[i]]) {
				++z[i];
			}
			
			if (i + z[i] - 1 > r) {
				l = i;
				r = i + z[i] - 1;
			}
		}
		
		return z;
	}
\end{lstlisting}

\complexity{Time: $O(n)$, memory: $O(n)$.}

\subsection*{Prefix Function}

\begin{lstlisting}[style=cpp]
	vector<ll> pref_func(string &s) {
		vector<ll> pref(s.size());
		
		for (int i = 1; i < (int)s.size(); i++) {
			ll k = pref[i - 1];
			
			while (k && s[k] != s[i]) {
				k = pref[k - 1];
			}
			
			if (s[k] == s[i]) {
				pref[i] = k + 1;
			}
		}
		
		return pref;
	}
\end{lstlisting}

\complexity{Time: $O(n)$, memory: $O(n)$.}

\subsection*{Aho-Corasick}

\begin{lstlisting}[style=cpp]
	struct aho_corasick {
		static const int alp = 26;
		
		struct node {
			node *tr[alp]{};
			node *suff_link = nullptr;
			node *to_pr = nullptr;
			int len = 0;
			vector<int> pr;
		};
		
		vector<unique_ptr<node>> nodes;
		node *root;
		bool built = false;
		
		node* new_node() {
			nodes.push_back(make_unique<node>());
			return nodes.back().get();
		}
		
		node* fnd(node *now, int need) {
			while (now != root && now->tr[need] == nullptr) {
				now = now->suff_link;
			}
			
			if (now->tr[need] != nullptr) {
				return now->tr[need];
			}
			
			return root;
		}
		
		public:
		aho_corasick() {
			root = new_node();
			root->suff_link = root;
			root->to_pr = root;
		}
		
		void add(const string &s, int id) {
			assert(!built);
			
			node *now = root;
			
			for (char ch : s) {
				int c = ch - 'a';
				assert(0 <= c && c < alp);
				
				if (now->tr[c] == nullptr) {
					now->tr[c] = new_node();
					now->tr[c]->len = now->len + 1;
				}
				
				now = now->tr[c];
			}
			
			now->pr.push_back(id);
		}
		
		void build() {
			assert(!built);
			
			queue<node*> q;
			
			for (int i = 0; i < alp; i++) {
				if (root->tr[i] != nullptr) {
					root->tr[i]->suff_link = root;
					root->tr[i]->to_pr = root;
					q.push(root->tr[i]);
				}
			}
			
			while (!q.empty()) {
				node *now = q.front();
				q.pop();
				
				for (int i = 0; i < alp; i++) {
					if (now->tr[i] == nullptr) {
						continue;
					}
					
					node *u = now->tr[i];
					
					u->suff_link = fnd(now->suff_link, i);
					
					if (!u->suff_link->pr.empty()) {
						u->to_pr = u->suff_link;
					}
					else {
						u->to_pr = u->suff_link->to_pr;
					}
					
					q.push(u);
				}
			}
			
			built = true;
		}
		
		vector<vector<int>> find(
		const string &text, int cnt
		) const {
			assert(built);
			
			vector<vector<int>> ans(cnt);
			node *now = root;
			
			for (int i = 0; i < (int)text.size(); i++) {
				int c = text[i] - 'a';
				assert(0 <= c && c < alp);
				
				while (now != root && now->tr[c] == nullptr) {
					now = now->suff_link;
				}
				
				if (now->tr[c] != nullptr) {
					now = now->tr[c];
				}
				else {
					now = root;
				}
				
				for (node *u = now; u != root; u = u->to_pr) {
					for (int id : u->pr) {
						ans[id].push_back(i - u->len + 1);
					}
				}
			}
			
			return ans;
		}
	};
\end{lstlisting}

\complexity{Добавление образцов: $O(S)$. Построение суффиксных ссылок: $O(VH)$ в худшем случае. Поиск: $O(T+Z)$. Память: $O(26V+Z)$.}

\subsection*{Fermat Primality Test}

\begin{lstlisting}[style=cpp]
	ll mult(ll a, ll b, ll mod) {
		return (__int128)a * b % mod;
	}
	
	ll bin_pow(ll a, ll n, ll mod) {
		ll res = 1;
		
		while (n > 0) {
			if (n % 2) {
				res = mult(res, a, mod);
			}
			
			n /= 2;
			a = mult(a, a, mod);
		}
		
		return res;
	}
	
	bool is_prime_fermat(ll n, int iterations = 50) {
		if (n < 2) return false;
		if (n <= 3) return true;
		if (n % 2 == 0) return false;
		
		// Trial division up to cube root.
		for (ll i = 3; (__int128)i * i * i <= n; i += 2) {
			if (n % i == 0) {
				return false;
			}
		}
		
		static mt19937_64 rng(
		chrono::steady_clock::now().time_since_epoch().count()
		);
		
		for (int i = 0; i < iterations; i++) {
			ll a = uniform_int_distribution<ll>(2, n - 2)(rng);
			
			if (std::gcd(a, n) != 1) {
				return false;
			}
			
			if (bin_pow(a, n - 1, n) != 1) {
				return false;
			}
		}
		
		return true;
	}
\end{lstlisting}

\complexity{Время: $O(n^{1/3} + k\log n)$, память: $O(1)$, где $k$ — количество итераций.}

\pitfall Вероятностный алгоритм: после предварительной проверки до кубического корня оставшиеся составные числа не являются числами Кармайкла. Вероятность ошибки при независимом равномерном выборе оснований не превосходит $2^{-k}$.

\subsection*{Implicit Treap}

\begin{lstlisting}[style=cpp]
	class Treap {
		static mt19937_64 generator;
		
		struct node {
			node *left = 0, *right = 0;
			ll value, mmin, sum, add;
			ll size, prio;
			
			node(ll v) : value(v), mmin(v), sum(v),
			size(1), prio(generator()), add(0LL) {}
		};
		
		node *head = 0;
		
		static ll get_size(node *ver) {
			return ver ? ver->size : 0;
		}
		
		static ll get_min(node *ver) {
			return ver ? ver->mmin : numeric_limits<ll>::max();
		}
		
		static ll get_sum(node *ver) {
			return ver ? ver->sum : 0LL;
		}
		
		static void push(node *ver) {
			if (!ver)
			return;
			
			if (ver->add) {
				ver->mmin += ver->add;
				ver->value += ver->add;
				ver->sum += ver->size * ver->add;
				
				if (ver->left)
				ver->left->add += ver->add;
				
				if (ver->right)
				ver->right->add += ver->add;
				
				ver->add = 0;
			}
		}
		
		static void update(node *ver) {
			if (ver) {
				push(ver->left);
				push(ver->right);
				
				ver->size = 1 + get_size(ver->left)
				+ get_size(ver->right);
				
				ver->mmin = min(ver->value,
				min(get_min(ver->left),
				get_min(ver->right)));
				
				ver->sum = ver->value
				+ get_sum(ver->left)
				+ get_sum(ver->right);
			}
		}
		
		static node* merge(node *a, node *b) {
			push(a);
			push(b);
			
			if (!a || !b)
			return a ? a : b;
			
			if (a->prio > b->prio) {
				a->right = merge(a->right, b);
				update(a);
				return a;
			}
			else {
				b->left = merge(a, b->left);
				update(b);
				return b;
			}
		}
		
		static void split(
		node *ver, int k, node *&a, node *&b
		) {
			push(ver);
			
			if (!ver) {
				a = 0;
				b = 0;
				return;
			}
			
			if (get_size(ver->left) < k) {
				split(ver->right,
				k - get_size(ver->left) - 1,
				ver->right, b);
				a = ver;
			}
			else {
				split(ver->left, k, a, ver->left);
				b = ver;
			}
			
			update(a);
			update(b);
		}
		
		public:
		void insert(int index, ll value) {
			node *less = 0, *greater = 0;
			
			split(head, index, less, greater);
			
			head = merge(
			merge(less, new node(value)),
			greater
			);
		}
		
		void add_range(int left, int right, ll new_value) {
			node *less = 0, *need = 0, *greater = 0;
			
			split(head, right + 1, less, greater);
			split(less, left, less, need);
			
			need->add += new_value;
			
			head = merge(merge(less, need), greater);
		}
		
		void swap_ranges(int l1, int r1, int l2, int r2) {
			node *a = 0, *b = 0;
			node *h1 = 0, *h2 = 0, *h3 = 0;
			
			split(head, r2 + 1, b, h3);
			split(b, l2, h2, b);
			split(h2, r1 + 1, a, h2);
			split(a, l1, h1, a);
			
			head = merge(h1, b);
			head = merge(head, h2);
			head = merge(head, a);
			head = merge(head, h3);
		}
		
		ll min_req(int left, int right) {
			node *less = 0, *greater = 0, *need = 0;
			
			split(head, right + 1, less, greater);
			split(less, left, less, need);
			
			ll ans = get_min(need);
			
			head = merge(merge(less, need), greater);
			
			return ans;
		}
		
		ll sum_req(int left, int right) {
			node *less = 0, *greater = 0, *need = 0;
			
			split(head, right + 1, less, greater);
			split(less, left, less, need);
			
			ll ans = get_sum(need);
			
			head = merge(merge(less, need), greater);
			
			return ans;
		}
	};
	
	mt19937_64 Treap::generator;
\end{lstlisting}

\complexity{Expected time per operation: $O(\log n)$. Memory: $O(n)$.}

\subsection*{Kuhn's Algorithm}

\begin{lstlisting}[style=cpp]
	struct kuhn {
		int n_left, n_right;
		int timer = 0;
		
		vector<vector<int>> graph1;
		vector<int> mt1, was1;
		
		kuhn(int n_left, int n_right)
		: n_left(n_left), n_right(n_right),
		graph1(n_left), mt1(n_right, -1),
		was1(n_left, 0) {}
		
		void add_edge(int u, int v) {
			graph1[u].push_back(v);
		}
		
		bool try_kuhn1(int ver) {
			if (was1[ver] == timer) {
				return false;
			}
			
			was1[ver] = timer;
			
			for (int to : graph1[ver]) {
				if (mt1[to] == -1 || try_kuhn1(mt1[to])) {
					mt1[to] = ver;
					return true;
				}
			}
			
			return false;
		}
		
		int max_matching() {
			fill(mt1.begin(), mt1.end(), -1);
			
			int ans = 0;
			
			for (int i = 0; i < n_left; i++) {
				++timer;
				
				if (try_kuhn1(i)) {
					ans++;
				}
			}
			
			return ans;
		}
	};
\end{lstlisting}

\complexity{Время: $O(L(E+1))$, память: $O(L+R+E)$.}

\subsection*{Fenwick Tree (1D)}

\begin{lstlisting}[style=cpp]
	struct fenwick {
		int n;
		vector<ll> F;
		
		fenwick(int n) : n(n), F(n + 1, 0) {}
		
		void update(int x, ll val = 1) {
			for (; x <= n; x += (x & -x)) {
				F[x] += val;
			}
		}
		
		// Sum on prefix [1, x].
		ll sum(int x) {
			ll s = 0;
			
			for (; x > 0; x -= (x & -x)) {
				s += F[x];
			}
			
			return s;
		}
		
		// Sum on segment [l, r].
		ll query(int l, int r) {
			return sum(r) - sum(l - 1);
		}
	};
\end{lstlisting}

\complexity{Update: $O(\log n)$, query: $O(\log n)$, memory: $O(n)$.}

\subsection*{Fenwick Tree (2D)}

\begin{lstlisting}[style=cpp]
	struct fenwick2d {
		int n, m;
		vector<vector<ll>> F;
		
		fenwick2d(int n, int m)
		: n(n), m(m), F(n + 1, vector<ll>(m + 1, 0)) {}
		
		void update(int x, int y, ll val = 1) {
			for (int i = x; i <= n; i += (i & -i)) {
				for (int j = y; j <= m; j += (j & -j)) {
					F[i][j] += val;
				}
			}
		}
		
		// Sum on rectangle [1, x] x [1, y].
		ll sum(int x, int y) {
			ll s = 0;
			
			for (int i = x; i > 0; i -= (i & -i)) {
				for (int j = y; j > 0; j -= (j & -j)) {
					s += F[i][j];
				}
			}
			
			return s;
		}
		
		// Sum on rectangle [x1, x2] x [y1, y2].
		ll query(int x1, int y1, int x2, int y2) {
			return sum(x2, y2)
			- sum(x1 - 1, y2)
			- sum(x2, y1 - 1)
			+ sum(x1 - 1, y1 - 1);
		}
	};
\end{lstlisting}

\complexity{Update: $O(\log n\log m)$, query: $O(\log n\log m)$, memory: $O(nm)$.}

\subsection*{Convex Polygon: Tangents and Point Inclusion}

\begin{lstlisting}[style=cpp]
	template<typename T>
	struct point {
		T x, y;
		
		point() : x(0), y(0) {}
		point(T x, T y) : x(x), y(y) {}
		
		point operator-(const point &other) const {
			return point(x - other.x, y - other.y);
		}
	};
	
	template<typename T>
	__int128 crs(const point<T> &a, const point<T> &b) {
		return (__int128)a.x * b.y -
		(__int128)a.y * b.x;
	}
	
	// Returns indices of two tangent vertices.
	// p must be strictly outside the polygon.
	pair<ll, ll> find_tangents(
	point<ll> p, const vector<point<ll>> &po
	) {
		ll n = po.size();
		assert(n >= 3);
		
		ll i, j, k, w;
		point<ll> a, b, c;
		
		auto magic = [&](ll coef) {
			k = 1;
			j = 0;
			
			while (k < n) {
				j++;
				k *= 2;
			}
			
			w = 0;
			
			for (i = j; i >= 0; --i) {
				a = po[(n + n + w - k) % n];
				b = po[w];
				c = po[(w + k) % n];
				
				if (crs(a - p, b - p) * coef >= 0) {
					if (crs(a - p, c - p) * coef >= 0) {
						w = (n + n + w - k) % n;
					}
					else {
						w = (w + k) % n;
					}
				}
				else {
					if (crs(c - p, b - p) * coef >= 0) {
						w = (w + k) % n;
					}
				}
				
				k /= 2;
			}
			
			return w;
		};
		
		pair<ll, ll> res;
		res.first = magic(1);
		res.second = magic(-1);
		
		return res;
	}
	
	// Checks whether p belongs to the convex polygon.
	// The boundary is considered inside.
	bool point_in_pol(
	point<ll> p, const vector<point<ll>> &po
	) {
		ll n = po.size();
		
		if (n < 3)
		return false;
		
		ll left = 1, right = n - 2, mid, pos;
		
		while (right - left > 1) {
			mid = (right + left) / 2;
			
			if (crs(po[mid] - po[0], p - po[0]) >= 0) {
				left = mid;
			}
			else {
				right = mid;
			}
		}
		
		if (crs(po[right] - po[0], p - po[0]) >= 0) {
			pos = right;
		}
		else {
			if (crs(po[left] - po[0], p - po[0]) >= 0) {
				pos = left;
			}
			else {
				return false;
			}
		}
		
		if (crs(p - po[0], po[pos + 1] - po[0]) >= 0 &&
		crs(po[pos + 1] - po[pos], p - po[pos]) >= 0) {
			return true;
		}
		
		return false;
	}
\end{lstlisting}

\complexity{Both algorithms: $O(\log n)$ time and $O(1)$ auxiliary memory.}


\subsection*{Primitive Root}

\begin{lstlisting}[style=cpp]
	ll primitive_root(ll m) {
		if (m < 2) return -1;
		if (m == 2) return 1;
		
		auto factorize = [](ll x) {
			vector<ll> res;
			
			for (ll i = 2; i <= x / i; i++) {
				while (x % i == 0) {
					res.push_back(i);
					x /= i;
				}
			}
			
			if (x > 1) res.push_back(x);
			
			return res;
		};
		
		auto has_primitive = [&](ll x) {
			if (x == 2 || x == 4) return true;
			
			vector<ll> s = factorize(x);
			
			if (s.size() == 1) return true;
			
			if (s[0] == 2) {
				if (s[1] == 2) return false;
				
				for (int i = 1; i < (int)s.size(); i++) {
					if (s[i] != s[1]) return false;
				}
				
				return true;
			}
			
			for (int i = 0; i < (int)s.size(); i++) {
				if (s[i] != s[0]) return false;
			}
			
			return true;
		};
		
		auto phi = [](ll x) {
			ll ans = 1;
			
			for (ll i = 2; i <= x / i; i++) {
				if (x % i == 0) {
					ll tt = i - 1;
					x /= i;
					
					while (x % i == 0) {
						tt *= i;
						x /= i;
					}
					
					ans *= tt;
				}
			}
			
			if (x > 1) ans *= (x - 1);
			
			return ans;
		};
		
		auto mult = [](ll a, ll b, ll mod) {
			return (__int128)a * b % mod;
		};
		
		auto bin_pow = [&](ll a, ll n, ll mod) {
			ll res = 1;
			
			while (n) {
				if (n & 1) {
					res = mult(res, a, mod);
				}
				
				a = mult(a, a, mod);
				n >>= 1;
			}
			
			return res;
		};
		
		auto gcd = [](ll a, ll b) {
			while (b) {
				ll t = a % b;
				a = b;
				b = t;
			}
			
			return a;
		};
		
		if (!has_primitive(m)) return -1;
		
		ll ph = phi(m);
		vector<ll> factors = factorize(ph);
		
		vector<ll> p;
		
		for (int i = 0; i < (int)factors.size(); i++) {
			if (i == 0 || factors[i] != factors[i - 1]) {
				p.push_back(ph / factors[i]);
			}
		}
		
		for (ll g = 2; g < m; g++) {
			if (gcd(g, m) != 1) continue;
			
			bool flag = true;
			
			for (ll c : p) {
				if (bin_pow(g, c, m) == 1) {
					flag = false;
					break;
				}
			}
			
			if (flag) return g;
		}
		
		return -1;
	}
\end{lstlisting}

\complexity{Время: $O(\sqrt m + G\cdot\omega(\varphi(m))\log m)$, память: $O(\log m)$, где $G$ — ответ, $\omega(x)$ — количество различных простых делителей $x$.}

Дан модуль $m$ и числа $a,b$. Необходимо найти минимальное
неотрицательное целое число $x$, удовлетворяющее сравнению
\[
a^x \equiv b \pmod m.
\]
Функция возвращает $-1$, если решения не существует.
Предполагается, что $\gcd(a,m)=1$.

Используется алгоритм Baby-step Giant-step (BSGS).

\complexity{Время: $O(\sqrt m\log m)$, память: $O(\sqrt m)$.}

\subsection*{Discrete Logarithm (BSGS)}

\begin{lstlisting}[style=cpp]
	ll discrete_log(ll a, ll b, ll mod) {
		if (mod == 1) return 0;
		
		assert(mod > 0);
		
		a %= mod;
		b %= mod;
		
		if (a < 0) a += mod;
		if (b < 0) b += mod;
		
		ll g = a, h = mod;
		while (h) {
			ll t = g % h;
			g = h;
			h = t;
		}
		
		if (g != 1) return -1;
		
		auto mul = [mod](ll x, ll y) {
			return (ll)((__int128)x * y % mod);
		};
		
		auto bin_pow = [&](ll x, ll n) {
			ll res = 1;
			
			while (n) {
				if (n & 1) res = mul(res, x);
				x = mul(x, x);
				n >>= 1;
			}
			
			return res;
		};
		
		auto get_inv = [&](ll x) {
			__int128 r0 = x, r1 = mod;
			__int128 x0 = 1, x1 = 0;
			
			while (r1) {
				__int128 q = r0 / r1;
				
				__int128 nr = r0 - q * r1;
				r0 = r1;
				r1 = nr;
				
				__int128 nx = x0 - q * x1;
				x0 = x1;
				x1 = nx;
			}
			
			if (r0 != 1) return -1LL;
			
			x0 %= mod;
			if (x0 < 0) x0 += mod;
			
			return (ll)x0;
		};
		
		ll s = sqrtl((long double)mod);
		while ((__int128)s * s < mod) s++;
		
		vector<pair<ll, ll>> baby;
		
		ll cur = 1;
		for (ll j = 0; j < s; j++) {
			baby.push_back({cur, j});
			cur = mul(cur, a);
		}
		
		sort(baby.begin(), baby.end());
		
		ll giant = get_inv(bin_pow(a, s));
		cur = b;
		
		for (ll i = 0; i < s; i++) {
			auto it = lower_bound(
			baby.begin(), baby.end(), make_pair(cur, -1LL)
			);
			
			if (it != baby.end() && it->first == cur) {
				return i * s + it->second;
			}
			
			cur = mul(cur, giant);
		}
		
		return -1;
	}
\end{lstlisting}

\complexity{Time: $O(\sqrt{mod}\log mod)$, memory: $O(\sqrt{mod})$.}

\subsection*{Modular Roots}

Даны простое число $p$, целые числа $a$ и $k$.
Необходимо найти все решения сравнения
\[
x^k \equiv a \pmod p.
\]
Возвращается вектор всех различных корней в отсортированном порядке.

\begin{lstlisting}[style=cpp]
	vector<ll> modular_roots(ll a, ll k, ll p) {
		assert(p >= 2 && k >= 1);
		
		a %= p;
		if (a < 0) a += p;
		
		if (a == 0) return {0};
		
		ll g = primitive_root(p);
		ll t = discrete_log(g, a, p);
		
		if (t == -1) return {};
		
		ll n = p - 1;
		ll d = std::gcd(k, n);
		
		if (t % d != 0) return {};
		
		ll kk = k / d;
		ll nn = n / d;
		ll tt = t / d;
		
		// Find inverse of kk modulo nn.
		ll x, y;
		gcd(kk, nn, x, y);
		
		ll y0 = 0;
		
		if (nn > 1) {
			x %= nn;
			if (x < 0) x += nn;
			
			y0 = (__int128)(tt % nn) * x % nn;
		}
		
		auto mul_mod = [p](ll a, ll b) {
			return (ll)((__int128)a * b % p);
		};
		
		auto bin_pow = [&](ll a, ll n) {
			ll res = 1;
			
			while (n) {
				if (n & 1) res = mul_mod(res, a);
				a = mul_mod(a, a);
				n >>= 1;
			}
			
			return res;
		};
		
		vector<ll> ans;
		
		for (ll i = 0; i < d; i++) {
			ll exp = y0 + i * nn;
			ans.push_back(bin_pow(g, exp));
		}
		
		sort(ans.begin(), ans.end());
		
		return ans;
	}
\end{lstlisting}

\complexity{Дополнительно к поиску первообразного корня: $O(\sqrt p\log p+d\log p)$, память: $O(\sqrt p+d)$, где $d=\gcd(k,p-1)$.}

\subsection*{Persistent Segment Tree}

Персистентное дерево отрезков позволяет хранить несколько
версий массива, используя общие неизменяемые поддеревья.
При обновлении создаются только новые вершины на пути
от корня до изменяемого элемента (Path Copying).

Поддерживаются точечное присваивание, запрос суммы на
отрезке и копирование произвольной версии массива.

\begin{lstlisting}[style=cpp]
	struct node{
		node *left, *right;
		ll sum;
		
		node(){
			left = nullptr;
			right = nullptr;
			sum = 0;
		}
	};
	
	struct per_segtree{
		vector<node*> roots;
		ll size;
		
		void build(vector<ll> &ar, node *now, ll lx, ll rx){
			if (rx - lx == 1){
				now->sum = ar[lx];
				return;
			}
			
			if (now->left == nullptr){
				now->left = new node();
			}
			
			if (now->right == nullptr){
				now->right = new node();
			}
			
			build(ar, now->left, lx, (lx + rx) / 2);
			build(ar, now->right, (lx + rx) / 2, rx);
			
			now->sum = now->left->sum + now->right->sum;
		}
		
		void build(vector<ll> &ar){
			size = ar.size();
			roots.push_back(new node());
			build(ar, roots.back(), 0, size);
		}
		
		ll get(ll l, ll r, node *now, ll lx, ll rx){
			if (r <= lx || rx <= l){
				return 0;
			}
			
			if (lx >= l && rx <= r){
				return now->sum;
			}
			
			return get(l, r, now->left, lx, (lx + rx) / 2)
			+ get(l, r, now->right, (lx + rx) / 2, rx);
		}
		
		ll get(ll l, ll r, ll k){
			return get(l, r, roots[k], 0, size);
		}
		
		node* set(ll i, ll v, node *now, ll lx, ll rx){
			node *res = new node();
			
			res->left = now->left;
			res->right = now->right;
			
			if (rx - lx == 1){
				res->sum = v;
				return res;
			}
			
			if (i < (lx + rx) / 2){
				res->left = set(i, v, now->left, lx, (lx + rx) / 2);
			}
			else{
				res->right = set(i, v, now->right, (lx + rx) / 2, rx);
			}
			
			res->sum = res->left->sum + res->right->sum;
			return res;
		}
		
		void set(ll i, ll v, ll k){
			roots[k] = set(i, v, roots[k], 0, size);
		}
		
		void add_copy(ll k){
			roots.push_back(roots[k]);
		}
	};
\end{lstlisting}

\complexity{Build: $O(n)$; point assignment: $O(\log n)$;
	range sum: $O(\log n)$; version copy: $O(1)$.}

\textbf{Memory:} $O(n+q\log n+c)$, where $q$ is the
number of updates and $c$ is the number of copied versions.

\subsection*{Random XOR Hashing}

\begin{lstlisting}[style=cpp]
	mt19937_64 rnd_gen(
	chrono::steady_clock::now().time_since_epoch().count()
	);
	
	using ull = unsigned long long;
	
	unordered_map<ll, ull> xor_values;
	
	ull get_xor_value(ll x) {
		auto it = xor_values.find(x);
		
		if (it != xor_values.end()) {
			return it->second;
		}
		
		return xor_values[x] = rnd_gen();
	}
\end{lstlisting}

\subsection*{Random Integer}

\begin{lstlisting}[style=cpp]
	ll random_int(ll l, ll r) {
		assert(l <= r);
		return uniform_int_distribution<ll>(l, r)(rnd_gen);
	}
\end{lstlisting}

\subsection*{Random Real}

\begin{lstlisting}[style=cpp]
	ld random_real(ld l, ld r) {
		assert(l <= r);
		
		if (l == r) return l;
		
		return uniform_real_distribution<ld>(l, r)(rnd_gen);
	}
\end{lstlisting}


\end{document}
